\chapter*{Abstract}

%The past few decades have witnessed a number of significant
%advances in laser technology, raising the achievable
%intensity by a few orders of magnitude beyond the
%relativistic limit and reducing pulse duration to the attosecond.
%Thanks top Mourou \footnote{who happens to be
%the Nobel prize 2018 for the developpement
%of the CPA technique, along with Strickland and ???}
%\cite{book:mourou_2001} 
%and his collaborators from the 
%Michigan laboratory, amongst others, 
%soon as 2001, Chirped Pulse Amplificated (CPA) lasers
%achieving intensities greater than $10^{21}\;{\rm Wcm}^{-2}$,
%were available to experimenters, as they previously did not have
%any better than $10^{15}\;{\rm Wcm^{-2}}$.
%It means researchers now have access to
%ultrahigh intensity applications allowing an even
%better probing of matter. Such field strengths make 
%significant relativistic, magnetic and alien effects that 
%one could just not afford to neglect anymore. 
%
%In this situation, theoretists have been struggling in the 
%2000's to keep up experiments because of either 
%the computational limitations and the extensive lengths 
%and cost a laboratory would have needed to go through to access 
%decent workstations, especially in most of Africa and other third 
%world countries, either the new found obsolescence 
%of the the routinely used tools and approximations. 
%The first problem has largely been tackled by the increasing 
%computational power available (thanks to the then valid Moore's law) 
%and the developpement of always more clever algorithms and approaches 
%to tackle the problem.  
%As for the second problem, it is exactly the point of this work. 
%An new formulation of the kernel of the Coulomb potential in the 
%momentum space is proposed, allowing to take advantage 
%of all the benefits of this space without paying the usual price. 
%
%Without any doubts, this formulation of the kernel will allow 
%to smoothly go to even higher intensities as it sets free 
%from quadratures and other numerical technics, well known 
%to quickly lose their accuracy with either increasing basis sizes.

  


